\newpage

\section{Einleitung}
    \section{Motivation und Problemstellung}
Obwohl die knotenbasierten Detektionsmechanismen der betrachteten Ladeinfrastruktur bereits eine beachtliche Erkennungsleistung aufweisen (siehe Abbildung~\ref{fig:auszug_threat_report}), besteht das Ziel, die Spezifität (True-Negative-Rate) sowie die Robustheit des Gesamtsystems weiter zu steigern. Lokale Anomalieerkennungsverfahren stoßen systembedingt an Grenzen, wenn netzwerkweite Veränderungen oder subtile, koordinierte Manipulationsversuche vorliegen \cite{lo2022egraphsage}, die auf Ebene einer einzelnen Wallbox nicht als bösartig identifiziert werden können.

Die Motivation dieser Arbeit liegt daher in der Entwicklung einer Methodik, die es einzelnen Ladepunkten ermöglicht, Kontextinformationen benachbarter Knoten in den Entscheidungsprozess einzubeziehen. Es wird vermutet, dass durch diesen kollaborativen Ansatz räumliche und zeitliche Korrelationen innerhalb des Netzwerks besser erfasst werden können, die über die isolierte Betrachtung eines einzelnen Knotens hinausgehen. Durch die Integration dieser Kontextinformationen wird erhofft, eine präzisere Differenzierung zwischen globalen Schwankungen und lokalen Sicherheitsvorfällen zu ermöglichen. Es wird erwartet, dass dieser erweiterte Kontextbezug zu einer Reduktion von Fehlalarmen (False Positives) sowie einer verbesserten Erkennungsrate tatsächlicher Bedrohungen (True Positives) beiträgt. Inwiefern diese Annahmen empirisch belastbar sind, bleibt jedoch offen und soll in einer weiterführenden Arbeit durch die Evaluation anhand eines unabhängigen Datensatzes systematisch untersucht werden.
    \section{Zielsetzung der Arbeit}
	Um die dargelegte Problemstellung zu adressieren, liegt der Fokus dieser Arbeit auf der Konzeption und Implementierung einer robusten Infrastruktur für die Anomalieerkennung innerhalb der Ladeinfrastruktur. Übergeordnetes Ziel ist die Schaffung einer stabilen Betriebsumgebung, die einen zuverlässigen, kontextbasierten Datenaustausch zwischen den einzelnen Knoten gewährleistet. 

	Ein wesentlicher Aspekt hierbei ist die Realisierung einer konsistenten und effizienten Aggregation dezentraler Informationen (siehe Kapitel~\ref{sec:generierung_der_threat_reports} (Generierung der Threat Reports)), um die Detektionsgüte der beteiligten Entitäten signifikant zu steigern. Nur durch eine performante Umgebung, welche den reibungslosen Austausch von Kontextinformationen sicherstellt, kann die in Kapitel \ref{chapter:systemarchitektur} (Systemarchitektur und Detektion Framework) detaillierte Anomalieerkennung erfolgreich implementiert und langfristig wartbar betrieben werden. Die Entwicklung dieser stabilen Kommunikations- und Verarbeitungsarchitektur bildet somit den zentralen Bestandteil der vorliegenden Arbeit.
    \section{Aufbau der Arbeit}

Die vorliegende Arbeit gliedert sich wie folgt: Zunächst wird in \textbf{Kapitel \ref{chapter:systemarchitektur}}  die \textbf{Systemarchitektur} des ReSiLENT-Systems dargelegt. Der Fokus liegt dabei auf der dezentralen Kommunikation via IPFS sowie den vier Detektionseinheiten, welche physikalische, systemnahe und netzwerkbasierte Anomalien auf Wallbox-Ebene erfassen.

\textbf{Kapitel \ref{chapter:thread_report}} beschreibt die \textbf{Konzeption des Threat Reporters}. Es wird erläutert, wie Detektionsergebnisse unter Anwendung zeitlicher Diskretisierung und Sliding-Window-Analysen in einem standardisierten CSV-Format aggregiert und im Peer-to-Peer-Netzwerk distribuiert werden.

In \textbf{Kapitel \ref{chapter:theoretische_grundlagen_der_informationsfusion}} werden die \textbf{theoretischen Grundlagen der Informationsfusion} adressiert. Angesichts variabler Eingabemengen wird die methodische Wahl einer Kombination aus \textit{Deep Sets} zur Gewährleistung der Permutationsinvarianz und \textit{Attention}-Mechanismen zur dynamischen Relevanzbewertung begründet.

Darauf aufbauend stellt \textbf{Kapitel \ref{chapter:methodik_gewichtete_agg_von_threat_reports}} die \textbf{Methodik der gewichteten Aggregation} vor. Hierbei werden die \textit{Negative Exponential Utility Function} sowie eine \textit{logistische Sigmoid-Funktion} hinsichtlich ihres asymptotischen Sättigungsverhaltens und ihrer Robustheit gegenüber statistischen Ausreißern evaluiert.

Die \textbf{Implementierung} und Integration dieser Modelle in den ReSiLENT-Datenfluss werden in \textbf{Kapitel \ref{chapter:evaluation}} beschrieben. Die resultierenden aggregierten CTI-Reports werden präsentiert und hinsichtlich ihrer Qualität für die knoteninterne Nutzung bewertet.

Abschließend erfolgt in \textbf{Kapitel \ref{chapter:diskussion_und_forschungsluecke}} eine \textbf{Diskussion der Ergebnisse}. Dabei werden kritische Forschungslücken identifiziert, insbesondere die empirische Validierung der Gewichtungsstrategien, das Potenzial LLM-gestützter Aufbereitung für SOC-Operatoren sowie die Untersuchung systemischer Risiken wie des \textit{Echo-Chamber-Effekts}.


\newpage

\section{Systemarchitektur und Detection Framework}
\label{sec:systemarchitektur}
    \subsection{Architektur des Multi-Node-Systems}
    % TODO: Hier das Bild der Architektur einfügen
    \subsection{Analyse der implementierten Detections}
        \subsubsection{IDS Power und Temperature Monitoring}
        \subsubsection{Weitere Kern-Detections (6er Set)}
    \subsection{Datenakquise und Feature-Extraktion via Prometheus}
        \subsubsection{Prometheus Exporter Design}
        \subsubsection{Metriken: Durchschnitt, Min, Max und Median}

\newpage

\section{Konzeption des Threat Reporters}
    \subsection{Standardisierung durch den Threat Report (CSV/STIX)}
    \subsection{Zeitfenster-Analyse und Datenaggregation (60s/30s Intervalle)}
    \subsection{Schnittstellendesign für den Datentransfer zwischen Nodes}

\newpage

\section{Theoretische Grundlagen der Informationsfusion}
    \subsection{Vergleich von Aggregationsmethoden}
        \subsubsection{Dempster-Shafer-Evidenztheorie}
        \subsubsection{Fuzzy-basierte Fusion}
        \subsubsection{Bayes-basierte Konsensmodelle}
        \subsubsection{Gewichtetes Voting}
    \subsection{State-of-the-Art: Deep Sets und Attention-Mechanismen}
    % TODO: Wissenschaftliche Begründung der Methodenwahl (Deep Sets + Attention) hier einarbeiten

\newpage

\section{Methodik: Gewichtete Aggregation mittels Utility Theory}
    \subsection{Anforderungen an die Gewichtungsfunktion}
    \subsection{Mathematische Modellierung}
        \subsubsection{Die Gewichtungsformel: $w(n) = \text{BASE\_WEIGHT} + (\text{MAX\_WEIGHT} - \text{BASE\_WEIGHT}) \times (1 - \exp(-K \cdot n))$}
        \subsubsection{Parameteranalyse: Monotonie und Diminishing Returns}
        \subsubsection{Abgrenzung zur originären Utility Theory}
        % TODO: Die Adaption der Utility Theory für Gewichtungsfunktionen methodisch begründen
        %       und ggf. alternative Ansätze abwägen; Einschränkungen in der Diskussion thematisieren
    \subsection{Implementierung des Attention-basierten Selektionsmodells}

\newpage

\section{Erweiterung der Systemarchitektur und Implementierung}
    \subsection{Generierte Threat Reports: Format und Struktur}
    \subsection{Integration von Threat Reporter und Aggregationslogik}
    \subsection{Architektur der Fusionskomponente im P2P-Netzwerk}
    % TODO Grafik einfügen wie die Architektur nun verändert wurde
    
\newpage

\section{Evaluation}
    \subsection{Qualität der aggregierten Threat Reports}
    \subsection{Nützlichkeit für die interne Nutzung der Knoten}

\newpage

\section{Diskussion und Forschungslücken}
    \subsection{Auswirkungen unterschiedlicher Gewichtungsstrategien auf die Report-Qualität}
    \subsection{Potenzial der LLM-gestützten Aufbereitung für externe CTI-Nutzung}
    \subsection{Untersuchung des Einflusses von Report-Aggregation auf Detectionsergebnisse (Echo-Chamber-Effekt?)}

\newpage

\section{Zusammenfassung und Fazit}

\newpage
